\documentclass[10pt,a4paper]{amsart}
\usepackage{amsmath,amssymb,amsthm,graphicx}
% Retained preamble for the Stacks Project chapter snapshot in this directory.
% The chapter's own preamble.tex is a copy of the shared Stacks Project preamble; the block
% below is sliced MECHANICALLY, line for line, from it. Nothing was added, removed or
% redefined; the wrapper class, the xr-hyper/hyperref wiring and the externaldocument list
% are omitted because the assembler supplies the document class.
%
% lines taken from preamble.tex: 3, 5, 9, 11, 12, 156, 157, 158, ... (51 lines total)
%
% package coverage restored from the chapter's own preamble (2026-10-13): xy, xy-2cell

\usepackage[all]{xy}
\xyoption{2cell}
\UseAllTwocells

\newenvironment{reference}{\comment}{\endcomment}
\newenvironment{slogan}{\comment}{\endcomment}
\newenvironment{history}{\comment}{\endcomment}
\theoremstyle{plain}
\newtheorem{theorem}[subsection]{Theorem}
\newtheorem{proposition}[subsection]{Proposition}
\newtheorem{lemma}[subsection]{Lemma}
\theoremstyle{definition}
\newtheorem{definition}[subsection]{Definition}
\newtheorem{example}[subsection]{Example}
\newtheorem{exercise}[subsection]{Exercise}
\newtheorem{situation}[subsection]{Situation}
\theoremstyle{remark}
\newtheorem{remark}[subsection]{Remark}
\newtheorem{remarks}[subsection]{Remarks}
\numberwithin{equation}{subsection}
\def\lim{\mathop{\mathrm{lim}}\nolimits}
\def\colim{\mathop{\mathrm{colim}}\nolimits}
\def\Spec{\mathop{\mathrm{Spec}}}
\def\Hom{\mathop{\mathrm{Hom}}\nolimits}
\def\Ext{\mathop{\mathrm{Ext}}\nolimits}
\def\SheafHom{\mathop{\mathcal{H}\!\mathit{om}}\nolimits}
\def\SheafExt{\mathop{\mathcal{E}\!\mathit{xt}}\nolimits}
\def\Sch{\mathit{Sch}}
\def\Mor{\mathop{\mathrm{Mor}}\nolimits}
\def\Ob{\mathop{\mathrm{Ob}}\nolimits}
\def\Sh{\mathop{\mathit{Sh}}\nolimits}
\def\NL{\mathop{N\!L}\nolimits}
\def\CH{\mathop{\mathrm{CH}}\nolimits}
\def\proetale{{pro\text{-}\acute{e}tale}}
\def\etale{{\acute{e}tale}}
\def\QCoh{\mathit{QCoh}}
\def\Ker{\mathop{\mathrm{Ker}}}
\def\Im{\mathop{\mathrm{Im}}}
\def\Coker{\mathop{\mathrm{Coker}}}
\def\Coim{\mathop{\mathrm{Coim}}}
\DeclareMathSymbol{\boxtimes}{\mathbin}{AMSa}{"02}
\def\QCohstack{\mathcal{QC}\!\mathit{oh}}
\def\Cohstack{\mathcal{C}\!\mathit{oh}}
\def\Spacesstack{\mathcal{S}\!\mathit{paces}}
\def\Quotfunctor{\mathrm{Quot}}
\def\Hilbfunctor{\mathrm{Hilb}}
\def\Curvesstack{\mathcal{C}\!\mathit{urves}}
\def\Polarizedstack{\mathcal{P}\!\mathit{olarized}}
\def\Complexesstack{\mathcal{C}\!\mathit{omplexes}}
\def\Pic{\mathop{\mathrm{Pic}}\nolimits}
\def\Picardstack{\mathcal{P}\!\mathit{ic}}
\def\Picardfunctor{\mathrm{Pic}}
\def\Deformationcategory{\mathcal{D}\!\mathit{ef}}
\usepackage{cleveref}
\crefname{originallemma}{original lemma}{original lemmas}
\Crefname{originallemma}{Original lemma}{Original lemmas}
\makeatletter
\newlabel{curves-lemma-characterize-contraction-to-stable}{{0E7Q}{source}{}{sourcepointer-1}{}}
\newlabel{curves-lemma-characterize-contraction-to-stable@cref}{{[originallemma][0][]0E7Q}{[1][0][]source}{}{}{}}
\makeatother
\crefname{originallemma}{original lemma}{original lemmas}
\Crefname{originallemma}{Original lemma}{Original lemmas}
\makeatletter
\newlabel{lemma-contract-basic-uniqueness}{{0E89}{source}{}{sourcepointer-2}{}}
\newlabel{lemma-contract-basic-uniqueness@cref}{{[originallemma][0][]0E89}{[1][0][]source}{}{}{}}
\makeatother
\crefname{originallemma}{original lemma}{original lemmas}
\Crefname{originallemma}{Original lemma}{Original lemmas}
\makeatletter
\newlabel{lemma-contract-prestable-to-stable}{{0E8A}{source}{}{sourcepointer-3}{}}
\newlabel{lemma-contract-prestable-to-stable@cref}{{[originallemma][0][]0E8A}{[1][0][]source}{}{}{}}
\makeatother
\crefname{originallemma}{original lemma}{original lemmas}
\Crefname{originallemma}{Original lemma}{Original lemmas}
\makeatletter
\newlabel{models-example-blowup}{{0CDC}{source}{}{sourcepointer-4}{}}
\newlabel{models-example-blowup@cref}{{[originallemma][0][]0CDC}{[1][0][]source}{}{}{}}
\makeatother
\crefname{originallemma}{original lemma}{original lemmas}
\Crefname{originallemma}{Original lemma}{Original lemmas}
\makeatletter
\newlabel{models-lemma-blowdown-at-worst-nodal}{{0CDF}{source}{}{sourcepointer-5}{}}
\newlabel{models-lemma-blowdown-at-worst-nodal@cref}{{[originallemma][0][]0CDF}{[1][0][]source}{}{}{}}
\makeatother
\crefname{originallemma}{original lemma}{original lemmas}
\Crefname{originallemma}{Original lemma}{Original lemmas}
\makeatletter
\newlabel{models-lemma-blowup-at-worst-nodal}{{0CDE}{source}{}{sourcepointer-6}{}}
\newlabel{models-lemma-blowup-at-worst-nodal@cref}{{[originallemma][0][]0CDE}{[1][0][]source}{}{}{}}
\makeatother
\crefname{originallemma}{original lemma}{original lemmas}
\Crefname{originallemma}{Original lemma}{Original lemmas}
\makeatletter
\newlabel{models-lemma-minimal-model-unique}{{0C6B}{source}{}{sourcepointer-7}{}}
\newlabel{models-lemma-minimal-model-unique@cref}{{[originallemma][0][]0C6B}{[1][0][]source}{}{}{}}
\makeatother
\crefname{originallemma}{original lemma}{original lemmas}
\Crefname{originallemma}{Original lemma}{Original lemmas}
\makeatletter
\newlabel{models-lemma-pre-exists-minimal-model}{{0CD9}{source}{}{sourcepointer-8}{}}
\newlabel{models-lemma-pre-exists-minimal-model@cref}{{[originallemma][0][]0CD9}{[1][0][]source}{}{}{}}
\makeatother
\crefname{originallemma}{original lemma}{original lemmas}
\Crefname{originallemma}{Original lemma}{Original lemmas}
\makeatletter
\newlabel{models-lemma-semistable}{{0CDG}{source}{}{sourcepointer-9}{}}
\newlabel{models-lemma-semistable@cref}{{[originallemma][0][]0CDG}{[1][0][]source}{}{}{}}
\makeatother
\crefname{originallemma}{original lemma}{original lemmas}
\Crefname{originallemma}{Original lemma}{Original lemmas}
\makeatletter
\newlabel{models-proposition-exists-minimal-model}{{0C2W}{source}{}{sourcepointer-10}{}}
\newlabel{models-proposition-exists-minimal-model@cref}{{[originallemma][0][]0C2W}{[1][0][]source}{}{}{}}
\makeatother
\crefname{originallemma}{original lemma}{original lemmas}
\Crefname{originallemma}{Original lemma}{Original lemmas}
\makeatletter
\newlabel{models-section-models}{{0C2R}{source}{}{sourcepointer-11}{}}
\newlabel{models-section-models@cref}{{[originallemma][0][]0C2R}{[1][0][]source}{}{}{}}
\makeatother
\crefname{originallemma}{original lemma}{original lemmas}
\Crefname{originallemma}{Original lemma}{Original lemmas}
\makeatletter
\newlabel{more-morphisms-lemma-check-h1-fibre-zero}{{0E7F}{source}{}{sourcepointer-12}{}}
\newlabel{more-morphisms-lemma-check-h1-fibre-zero@cref}{{[originallemma][0][]0E7F}{[1][0][]source}{}{}{}}
\makeatother
\crefname{originallemma}{original lemma}{original lemmas}
\Crefname{originallemma}{Original lemma}{Original lemmas}
\makeatletter
\newlabel{more-morphisms-lemma-h1-fibre-zero}{{0E7G}{source}{}{sourcepointer-13}{}}
\newlabel{more-morphisms-lemma-h1-fibre-zero@cref}{{[originallemma][0][]0E7G}{[1][0][]source}{}{}{}}
\makeatother
\crefname{originallemma}{original lemma}{original lemmas}
\Crefname{originallemma}{Original lemma}{Original lemmas}
\makeatletter
\newlabel{more-morphisms-lemma-h1-fibre-zero-check-h0-kappa}{{0E7I}{source}{}{sourcepointer-14}{}}
\newlabel{more-morphisms-lemma-h1-fibre-zero-check-h0-kappa@cref}{{[originallemma][0][]0E7I}{[1][0][]source}{}{}{}}
\makeatother
\crefname{originallemma}{original lemma}{original lemmas}
\Crefname{originallemma}{Original lemma}{Original lemmas}
\makeatletter
\newlabel{resolve-lemma-cohomology-of-blowup}{{0AGS}{source}{}{sourcepointer-15}{}}
\newlabel{resolve-lemma-cohomology-of-blowup@cref}{{[originallemma][0][]0AGS}{[1][0][]source}{}{}{}}
\makeatother
\crefname{originallemma}{original lemma}{original lemmas}
\Crefname{originallemma}{Original lemma}{Original lemmas}
\makeatletter
\newlabel{resolve-section-minus-one}{{0C2I}{source}{}{sourcepointer-16}{}}
\newlabel{resolve-section-minus-one@cref}{{[originallemma][0][]0C2I}{[1][0][]source}{}{}{}}
\makeatother
\title{Stacks, Moduli of Curves: lemma-unique-stable-model}
\author{The Stacks Project Authors}
\date{stacks-project:moduli-curves.tex:local-snapshot-20261012; GNU FDL 1.2 or later}
\begin{document}
\maketitle
\noindent\textit{Editorial note.} Statement span is lines 3272-3281 of sources/stacks-moduli-curves/pinned-original-source.tex; proof span is lines 3283-3376, the complete author proof (92 body lines).\par\medskip
\noindent\textit{Editorial note.} Navigation stubs: 16 cited-but-not-included labels are supplied as non-load-bearing pointers with their pinned tags where the snapshot tag index has them.\par\medskip
\noindent\textit{Editorial note.} extra\_packages is empty on purpose: the retained preamble already loads xy with options and a second declaration clashes in the assembler wrapper.\par\medskip
\noindent\textit{Editorial note.} Text is retained verbatim, including the author's own typos and placeholders.\par\medskip
\noindent\textit{Editorial note.} Queue index 2 of queue\_moduli-curves.json (pb=4146, stmt 3272-3281, proof 3283-3376).\par\medskip
\noindent\textit{Allowed standard prerequisites (editorial).}\begin{itemize}
\item Explicit prerequisite boundary. The cited inputs are prerequisites; the argument in the retained span is not replaced by any cited result.
\end{itemize}\medskip
\ifdefined\hypertarget\hypertarget{sourcepointer-1}{}\fi
\noindent\textit{Source reference.} Original lemma 0E7Q: sources/stacks-moduli-curves/pinned-original-source.tex, source lines 3376--3376. Cited by the retained statement or proof but not part of the collected spans; supplied only as a navigation pointer. The argument does not depend on its content.\par\medskip
\ifdefined\hypertarget\hypertarget{sourcepointer-2}{}\fi
\noindent\textit{Source reference.} Original lemma 0E89: sources/stacks-moduli-curves/pinned-original-source.tex, source lines 2900--2912. Cited by the retained statement or proof but not part of the collected spans; supplied only as a navigation pointer. The argument does not depend on its content.\par\medskip
\ifdefined\hypertarget\hypertarget{sourcepointer-3}{}\fi
\noindent\textit{Source reference.} Original lemma 0E8A: sources/stacks-moduli-curves/pinned-original-source.tex, source lines 3127--3138. Cited by the retained statement or proof but not part of the collected spans; supplied only as a navigation pointer. The argument does not depend on its content.\par\medskip
\ifdefined\hypertarget\hypertarget{sourcepointer-4}{}\fi
\noindent\textit{Source reference.} Original lemma 0CDC: sources/stacks-moduli-curves/pinned-original-source.tex, source lines 3376--3376. Cited by the retained statement or proof but not part of the collected spans; supplied only as a navigation pointer. The argument does not depend on its content.\par\medskip
\ifdefined\hypertarget\hypertarget{sourcepointer-5}{}\fi
\noindent\textit{Source reference.} Original lemma 0CDF: sources/stacks-moduli-curves/pinned-original-source.tex, source lines 3376--3376. Cited by the retained statement or proof but not part of the collected spans; supplied only as a navigation pointer. The argument does not depend on its content.\par\medskip
\ifdefined\hypertarget\hypertarget{sourcepointer-6}{}\fi
\noindent\textit{Source reference.} Original lemma 0CDE: sources/stacks-moduli-curves/pinned-original-source.tex, source lines 3376--3376. Cited by the retained statement or proof but not part of the collected spans; supplied only as a navigation pointer. The argument does not depend on its content.\par\medskip
\ifdefined\hypertarget\hypertarget{sourcepointer-7}{}\fi
\noindent\textit{Source reference.} Original lemma 0C6B: sources/stacks-moduli-curves/pinned-original-source.tex, source lines 3376--3376. Cited by the retained statement or proof but not part of the collected spans; supplied only as a navigation pointer. The argument does not depend on its content.\par\medskip
\ifdefined\hypertarget\hypertarget{sourcepointer-8}{}\fi
\noindent\textit{Source reference.} Original lemma 0CD9: sources/stacks-moduli-curves/pinned-original-source.tex, source lines 3376--3376. Cited by the retained statement or proof but not part of the collected spans; supplied only as a navigation pointer. The argument does not depend on its content.\par\medskip
\ifdefined\hypertarget\hypertarget{sourcepointer-9}{}\fi
\noindent\textit{Source reference.} Original lemma 0CDG: sources/stacks-moduli-curves/pinned-original-source.tex, source lines 3376--3376. Cited by the retained statement or proof but not part of the collected spans; supplied only as a navigation pointer. The argument does not depend on its content.\par\medskip
\ifdefined\hypertarget\hypertarget{sourcepointer-10}{}\fi
\noindent\textit{Source reference.} Original lemma 0C2W: sources/stacks-moduli-curves/pinned-original-source.tex, source lines 3376--3376. Cited by the retained statement or proof but not part of the collected spans; supplied only as a navigation pointer. The argument does not depend on its content.\par\medskip
\ifdefined\hypertarget\hypertarget{sourcepointer-11}{}\fi
\noindent\textit{Source reference.} Original lemma 0C2R: sources/stacks-moduli-curves/pinned-original-source.tex, source lines 3376--3376. Cited by the retained statement or proof but not part of the collected spans; supplied only as a navigation pointer. The argument does not depend on its content.\par\medskip
\ifdefined\hypertarget\hypertarget{sourcepointer-12}{}\fi
\noindent\textit{Source reference.} Original lemma 0E7F: sources/stacks-moduli-curves/pinned-original-source.tex, source lines 3376--3376. Cited by the retained statement or proof but not part of the collected spans; supplied only as a navigation pointer. The argument does not depend on its content.\par\medskip
\ifdefined\hypertarget\hypertarget{sourcepointer-13}{}\fi
\noindent\textit{Source reference.} Original lemma 0E7G: sources/stacks-moduli-curves/pinned-original-source.tex, source lines 3376--3376. Cited by the retained statement or proof but not part of the collected spans; supplied only as a navigation pointer. The argument does not depend on its content.\par\medskip
\ifdefined\hypertarget\hypertarget{sourcepointer-14}{}\fi
\noindent\textit{Source reference.} Original lemma 0E7I: sources/stacks-moduli-curves/pinned-original-source.tex, source lines 3376--3376. Cited by the retained statement or proof but not part of the collected spans; supplied only as a navigation pointer. The argument does not depend on its content.\par\medskip
\ifdefined\hypertarget\hypertarget{sourcepointer-15}{}\fi
\noindent\textit{Source reference.} Original lemma 0AGS: sources/stacks-moduli-curves/pinned-original-source.tex, source lines 3376--3376. Cited by the retained statement or proof but not part of the collected spans; supplied only as a navigation pointer. The argument does not depend on its content.\par\medskip
\ifdefined\hypertarget\hypertarget{sourcepointer-16}{}\fi
\noindent\textit{Source reference.} Original lemma 0C2I: sources/stacks-moduli-curves/pinned-original-source.tex, source lines 3376--3376. Cited by the retained statement or proof but not part of the collected spans; supplied only as a navigation pointer. The argument does not depend on its content.\par\medskip
\begin{lemma}
\label{lemma-unique-stable-model}
Let $R$ be a discrete valuation ring with fraction field $K$.
Let $C$ be a smooth proper curve over $K$
with $K = H^0(C, \mathcal{O}_C)$ and genus $g$.
If $X$ and $X'$ are models of $C$
(Semistable Reduction, Section \ref{models-section-models})
and $X$ and $X'$ are stable families of genus $g$ curves over $R$,
then there exists a unique isomorphism $X \to X'$ of models.
\end{lemma}

\begin{proof}
Let $Y$ be the minimal model for $C$. Recall that $Y$ exists, is
unique, and is at-worst-nodal of relative dimension $1$ over $R$, see
Semistable Reduction,
Proposition \ref{models-proposition-exists-minimal-model} and
Lemmas \ref{models-lemma-minimal-model-unique} and
\ref{models-lemma-semistable} (applies because we have $X$).
There is a contraction morphism
$$
Y \longrightarrow Z
$$
such that $Z$ is a stable family of curves of genus $g$ over $R$
(Lemma \ref{lemma-contract-prestable-to-stable}). We claim
there is a unique isomorphism of models $X \to Z$.
By symmetry the same is true for $X'$ and this will finish the proof.

\medskip\noindent
By Semistable Reduction, Lemma \ref{models-lemma-blowup-at-worst-nodal}
there exists a sequence
$$
X_m \to \ldots \to X_1 \to X_0 = X
$$
such that $X_{i + 1} \to X_i$ is the blowing up of a closed point
$x_i$ where $X_i$ is singular, $X_i \to \Spec(R)$ is at-worst-nodal
of relative dimension $1$, and $X_m$ is regular.
By Semistable Reduction, Lemma \ref{models-lemma-pre-exists-minimal-model}
there is a sequence
$$
X_m = Y_n \to Y_{n - 1} \to \ldots \to Y_1 \to Y_0 = Y
$$
of proper regular models of $C$, such that each morphism is a
contraction of an exceptional curve of the first kind\footnote{In fact
we have $X_m = Y$, i.e., $X_m$ does not contain any exceptional curves
of the first kind. We encourage the reader to think this through
as it simplifies the proof somewhat.}.
By Semistable Reduction, Lemma \ref{models-lemma-blowdown-at-worst-nodal}
each $Y_i$ is at-worst-nodal of relative dimension $1$ over $R$.
To prove the claim it suffices to show that there is an isomorphism
$X \to Z$ compatible with the morphisms $X_m \to X$
and $X_m = Y_n \to Y \to Z$. Let $s \in \Spec(R)$ be the closed point.
By either
Lemma \ref{lemma-contract-basic-uniqueness} or
Lemma \ref{lemma-contract-prestable-to-stable}
we reduce to proving that the morphisms
$X_{m, s} \to X_s$ and
$X_{m, s} \to Z_s$
are both equal to the canonical morphism of
Algebraic Curves, Lemma \ref{curves-lemma-characterize-contraction-to-stable}.

\medskip\noindent
For a morphism $c : U \to V$ of schemes over $\kappa(s)$
we say $c$ has property (*) if $\dim(U_v) \leq 1$ for $v \in V$,
$\mathcal{O}_V = c_*\mathcal{O}_U$, and $R^1c_*\mathcal{O}_U = 0$.
This property is stable under composition.
Since both $X_s$ and $Z_s$ are stable genus $g$ curves over $\kappa(s)$,
it suffices to show that each of the morphisms $Y_s \to Z_s$,
$X_{i + 1, s} \to X_{i, s}$, and $Y_{i + 1, s} \to Y_{i, s}$,
satisfy property (*), see
Algebraic Curves, Lemma \ref{curves-lemma-characterize-contraction-to-stable}.

\medskip\noindent
Property (*) holds for $Y_s \to Z_s$ by construction.

\medskip\noindent
The morphisms $c : X_{i + 1, s} \to X_{i, s}$ are constructed and studied
in the proof of
Semistable Reduction, Lemma \ref{models-lemma-blowup-at-worst-nodal}.
It suffices to check (*) \'etale locally on $X_{i, s}$.
Hence it suffices to check (*) for the base change of the morphism
``$X_1 \to X_0$'' in Semistable Reduction, Example \ref{models-example-blowup}
to $R/\pi R$.
We leave the explicit calculation to the reader.

\medskip\noindent
The morphism $c : Y_{i + 1, s} \to Y_{i, s}$ is the restriction
of the blow down of an exceptional curve $E \subset Y_{i + 1}$
of the first kind, i.e., $b : Y_{i + 1} \to Y_i$ is a contraction of $E$,
i.e., $b$ is a blowing up of a regular point on the surface $Y_i$
(Resolution of Surfaces, Section \ref{resolve-section-minus-one}).
Then $\mathcal{O}_{Y_i} = b_*\mathcal{O}_{Y_{i + 1}}$ and
$R^1b_*\mathcal{O}_{Y_{i + 1}} = 0$, see for example
Resolution of Surfaces, Lemma
\ref{resolve-lemma-cohomology-of-blowup}.
We conclude that $\mathcal{O}_{Y_{i, s}} = c_*\mathcal{O}_{Y_{i + 1, s}}$
and $R^1c_*\mathcal{O}_{Y_{i + 1, s}} = 0$ by
More on Morphisms, Lemmas \ref{more-morphisms-lemma-check-h1-fibre-zero},
\ref{more-morphisms-lemma-h1-fibre-zero}, and
\ref{more-morphisms-lemma-h1-fibre-zero-check-h0-kappa}
(only gives surjectivity of
$\mathcal{O}_{Y_{i, s}} \to c_*\mathcal{O}_{Y_{i + 1, s}}$ but
injectivity follows easily from the fact that $Y_{i, s}$ is reduced
and $c$ changes things only over one closed point).
This finishes the proof.
\end{proof}

\end{document}
